\section{Introduction}
\label{sec:introduction}

Legged robots can cross broken or narrow terrain as long as a sequence
of usable footholds exists. Finding that sequence requires more than
choosing foothold positions. The robot must also decide which leg
moves, when it lifts off, and how long its swing lasts. We call this
joint discrete decision \textit{contact selection}. We use
\textit{acyclic} for contact sequences that are neither a periodic
gait cycle nor transitions through a fixed gait library; the
literature also uses \textit{non-gaited} and \textit{free-gait}.

Most deployed model-based stacks fix the contact schedule in advance
and optimize motion around it \cite{Kim2019highly-dynamic,
Wensing2022Nov, grandia_perceptive_2022}. This is effective on regular
terrain, but a scheduled leg must step even when no good foothold is
reachable, and an unscheduled leg cannot step even when one is.

GaitNet makes only the contact decision. At each tick it outputs a
leg, a foothold, and a swing duration, or no action. It does
\emph{not} generate joint torques, joint positions, contact forces,
base trajectories, or whole-body control (WBC) objectives. A
downstream model-based controller realizes the decision. This
division keeps the learned component small enough to re-evaluate at a
high rate, and it keeps the foothold explicit so that the controller
can still check it against the terrain before execution.

\subsection{Related Work}
\label{subsec:related-work}

\textit{Optimization and search.} When contacts are decision
variables, planning becomes a mixed-integer program
\cite{Geisert2019contact-planning}. Convex terrain regions
\cite{deits_footstep_2014} and joint optimization of schedule and
timing \cite{winkler_gait_2018, aceituno_simultaneous_2019} shrink the
problem, but remain costly for fast reactive correction. Free-gait and
utility-based planners choose contact timing online from per-leg rules
\cite{fankhauser_free_gait_2016, boussema_online_2019,
bjelonic_whole-body_2021}, and search-based planners generate acyclic
sequences directly \cite{amatucci_monte_2022}. Recent work accelerates
search by learning cost-to-go \cite{taouil_non-gaited_2024} or
transition feasibility \cite{akizhanov_learning_2024}.
Contact-implicit MPC discovers contacts from dynamics in real time
\cite{kim_contact-implicit_2025}, but does not select footholds from a
perceived steppability map.

\textit{Learned controllers and hybrid planners.} End-to-end policies
cross difficult terrain reliably \cite{zhang_learning_2023,
zhang_novel_2024, duan_sim--real_2022, siekmann_blind_2021,
lee_learning_2020}, but the contact decision is implicit in joint
targets and cannot be inspected or constrained before execution.
Hybrid methods learn foothold placement while a gait generator
provides the schedule \cite{gangapurwala_rloc_2022,
villarreal_fast_2019, omar_fast_2022, omar_safesteps_2023,
asselmeier_steppability-informed_2024, gaspard_footstepnet_2024,
tsounis_deepgait_2020, shi_terrain-aware_2023, xie_glide_2023}, or
learn contact structure while a heuristic provides footholds
\cite{da_learning_2020, yang_fast_2021, sun_online_2024}. Patrizi et
al. \cite{patrizi_rl-augmented_2026} learn per-foot flight-phase
injections into an MPC horizon, which yields acyclic timing but not
foothold selection.

\textit{ContactNet.} The closest method is ContactNet
\cite{bratta_contactnet_2024}, which ranks discretized foothold
candidates with a neural network. Its concurrency is set by
configuration: one leg at a time in walking experiments and joint
footholds for a swinging pair in trot experiments. It re-ranks at
touchdown, and the configured gait sets swing timing.
\autoref{tab:capabilities} summarizes the differences we rely on.
Entries for ContactNet reflect only what is reported in
\cite{bratta_contactnet_2024}.

\begin{table}[t]
  \centering
  \caption{Contact-decision capabilities.}
  \label{tab:capabilities}
  \footnotesize
  \setlength{\tabcolsep}{2.5pt}
  \begin{tabular}{lccc}
    \toprule
    \textbf{Capability} & \textbf{Scheduled} & \textbf{ContactNet
    \cite{bratta_contactnet_2024}} & \textbf{GaitNet} \\
    \midrule
    Explicit foothold & Yes & Yes & Yes \\
    Chooses leg & Scheduled & Config.-dependent & Yes \\
    Chooses swing duration & Fixed & No & Yes$^{\dagger}$ \\
    Explicit no-action option & Implicit & \verify{Not reported} & Yes \\
    Re-decides during a swing & No & No (at touchdown) & Yes \\
    Swing overlap & Scheduled & Configured & Emergent \\
    \bottomrule
  \end{tabular}\\[2pt]
  \footnotesize{$^{\dagger}$The Go1 deployment floors the duration
  (\autoref{subsec:exp-hardware}).}
\end{table}

\subsection{Contributions}

\begin{enumerate}
  \item \textbf{High-rate acyclic selection.} A planner that selects
    one leg-specific foothold, swing duration, or no-action every
    \corlControlTickMs\,ms, without representing gait phase,
    predefined leg pairs, or a contact schedule.
  \item \textbf{Set-size-independent scoring.} A per-candidate
    architecture trained with sampled candidates but deployable
    through exhaustive batched enumeration without retraining.
  \item \textbf{Controller-decoupled execution.} A narrow
    planner-controller interface that lets the frozen planner operate
    with a model-based NMPC/WBC execution stack without accessing its
    internal optimization.
  \item \textbf{Systematic evaluation.} Matched Gazebo experiments
    studying terrain difficulty, concurrency restrictions, candidate
    enumeration, timing behavior, and computation, supplemented by a
    proof-of-concept Go1 deployment.
\end{enumerate}

An earlier thesis \cite{Sullivan2025} documented the initial
development of GaitNet. This paper provides its first archival
presentation and systematic evaluation, including exhaustive
deployment-time enumeration, controller transfer, concurrency
ablations, and physical deployment.
